\section{Exact and uniform approximate convexification of a pool block}
\label{app:convexification}

This appendix proves the convexification guarantees used in
Section~\ref{s6:synthesis}. The set is an isolated matrix block; no
physical quality or interpool constraints are retained. Throughout,
$N,m$ are positive integers and all error radii are nonnegative. Define
\[
 R_N=\{W\in\R_+^{N\times N}:\rank W\le1,
 W\mathbf1\le\mathbf1,\ W^T\mathbf1\le\mathbf1\},\qquad
 K_N=\conv(R_N).
\]
Both sets are compact: all entries of $R_N$ lie in $[0,1]$, and its
rank condition is closed. Every point of $K_N$ is a finite convex
combination of generators in $R_N$. The norm $\|W\|_1$ below is the
entrywise sum $\sum_{ij}|W_{ij}|$, with unit ball $U_1$; the dual
coefficient norm is $\|H\|_\infty=\max_{ij}|H_{ij}|$.

A lift over a cone $\mathcal C$ means an expression
$K=\pi(\mathcal C\cap L)$, where $L$ is affine and $\pi$ is linear
(or affine). Affine sections and images preserve the cone: intersect
$L$ with the preimage of the section equations and compose $\pi$ with
the desired map. Every scalar inequality counts as a nonnegative cone
factor. Semidefinite size means total matrix order, including scalar
blocks of order one. Arbitrary real formulation coefficients are allowed.

\subsection{An explicit correlation-polytope face}

\begin{theorem}[Correlation face]\label{app:correlation-face}
The affine section
\[
 F_m=\{W\in K_{2m}:\operatorname{tr}W=1,
 \ \sum_{i=1}^mW_{i,m+i}=0,\ \mathbf1^TW\mathbf1=m\}
\]
is a face affinely isomorphic to
$\operatorname{COR}(m)=\conv\{xx^T:x\in\{0,1\}^m\}$.
The map is $X=mW_{[m],[m]}$; writing $x=\operatorname{diag}X$, its
inverse is
\begin{equation}
 W=\frac1m
 \begin{pmatrix}
 X&x\mathbf1^T-X\\
 \mathbf1x^T-X&\mathbf1\mathbf1^T-x\mathbf1^T-\mathbf1x^T+X
 \end{pmatrix}.
 \label{app:face-inverse}
\end{equation}
\end{theorem}
\begin{proof}
For a nonzero generator $W=rc^T/S$, with $r,c\in[0,1]^{2m}$ of
common total $S$, one has $\operatorname{tr}W=r^Tc/S\le1$.
Equality forces $\sum_i r_i(1-c_i)=0$ and also
$\sum_i c_i(1-r_i)=0$. Hence $r_i>0$ implies $c_i=1$ and
$c_i>0$ implies $r_i=1$: both margins are the same nonempty binary
vector $a=\mathbf1_A$. The trace-one face is therefore the convex
hull of $aa^T/|A|$. Indeed, every positive-weight generator in a
convex combination attaining trace one must attain trace one itself.
The zero matrix cannot occur in that combination.

Nonnegativity makes the cross-pair equation another face equation.
It retains exactly those generators whose support contains at most
one element from each pair $\{i,m+i\}$. They satisfy $|A|\le m$.
Their total matrix mass is $|A|$, so the last equation selects the
maximal possible value $m$, exactly one element of every pair.
These are the atoms
\[
 \frac1m(x,\mathbf1-x)(x,\mathbf1-x)^T,
 \qquad x\in\{0,1\}^m.
\]
Successive faces are faces of the original set. The principal-block
map sends these atoms to $xx^T$, and block expansion gives
\eqref{app:face-inverse}. Its entries are affine in $X$, so the inverse
identity extends to the two convex hulls.
\end{proof}

\begin{corollary}[Exact conic size]\label{app:exact-conic}
For each fixed $d\ge1$, an exact lift of $K_{2m}$ over a product of
PSD cones of order at most $d$ needs $2^{\Omega(m)}$ blocks, with
constants depending on $d$. Every exact second-order-cone lift has
$2^{\Omega(m)}$ total cone dimension. Every exact SDP lift has total
matrix order $2^{\Omega(m^{2/13})}$. For $N\ge2$, $K_N$ has no finite
LP lift.
\end{corollary}
\begin{proof}
The face and affine map transfer any lift to $\operatorname{COR}(m)$
without changing its cone. The fixed-order bound follows from
\citet[Theorem~1]{s6:fawzi2013}; smaller blocks may be padded with
zeros to order $d$. A Lorentz cone of dimension $h\ge3$ can be
represented with $h-2$ three-dimensional Lorentz cones: introduce
successive partial Euclidean norms. A three-dimensional Lorentz cone
is linearly isomorphic to the cone of $2\times2$ PSD matrices.
Dimensions one and two use at most two nonnegative scalar factors.
Thus total Lorentz dimension $D$ gives at most $D$ blocks of order
two. The fixed-order lower bound proves the second assertion.
For unrestricted SDP order apply
\citet[Theorem~1.1]{s6:lrs-full}; several blocks combine as a
block-diagonal PSD matrix of their total order.

For nonpolyhedrality, consider $K_2$ on the face where its first row
and column margins equal one. Its generating matrices are
\[
 W(t)=\frac1{1+t}\begin{pmatrix}1&t\\t&t^2\end{pmatrix},
 \qquad 0\le t\le1.
\]
If $q=W_{11}=1/(1+t)$, then $W_{22}=q+1/q-2$ for
$1/2\le q\le1$. This function is strictly convex, so every graph
point is extreme in its convex hull: equality in the convexity
inequality for a convex combination requires all first coordinates
to agree. Therefore this face has infinitely many extreme points.
For $N>2$, fixing all entries outside a principal $2\times2$ block
to zero gives the same face. A projection of a finite polyhedron is
a polyhedron, so no finite LP lift is possible. The LP block-count
bound above is consequently only formal; the positive-error LP result
below is the substantive finite-size statement.
\end{proof}

\subsection{A quantitative transfer to the face}

For $W\in K_{2m}$ write
\[
 T(W)=\mathbf1^TW\mathbf1,\quad D(W)=1-\operatorname{tr}W,
 \quad B(W)=\sum_{i=1}^mW_{i,m+i}.
\]
The letter $D$ in this appendix denotes trace deficit, not the
interaction-cost matrix of Appendix~\ref{app:costs}.

\begin{lemma}[Exposure and rounding]\label{app:face-rounding}
Every $W\in K_{2m}$ satisfies
\begin{equation}
 T(W)\le m+2mB(W)+4mD(W).\label{app:exposure}
\end{equation}
There is $V\in F_m$ such that
\begin{equation}
 \|W-V\|_1\le28mB(W)+156mD(W)+5[m-T(W)].
 \label{app:face-distance}
\end{equation}
The right side is nonnegative; its last bracket need not be.
\end{lemma}
\begin{proof}
First take a nonzero generator $W=rc^T/S$. For $a,b\in[0,1]$,
$|a-b|\le a(1-b)+b(1-a)$, so
\begin{equation}
 \|r-c\|_1\le2SD(W).\label{app:margin-distance}
\end{equation}
The pair inequality $r_i+r_{m+i}-r_ir_{m+i}\le1$ yields
\[
 S\le m+\sum_i r_ir_{m+i}
 \le m+\sum_i r_ic_{m+i}+\|r-c\|_1
 \le m+SB+2SD.
\]
As $S\le2m$, this proves \eqref{app:exposure} for generators,
including zero directly, and hence for their convex hull by affinity.

Round each $r_i$ to a nearest binary value $a_i$. Since
$\min(t,1-t)\le2t(1-t)$ and
$\sum_i r_i(1-r_i)\le SD+\|r-c\|_1\le3SD$, we obtain
\[
 u:=\|r-a\|_1\le6SD,\qquad v:=\|c-a\|_1\le8SD.
\]
Let $h,z$ be the numbers of pairs containing respectively two and zero
ones in $a$, and put $k=\sum_i a_i=m+h-z$. The elementary bound
$|ab-cd|\le|a-c|+|b-d|$ in the unit box gives
\[
 h\le\sum_i r_ir_{m+i}+u
 \le SB+\|r-c\|_1+u\le SB+8SD.
\]
Change $a$ into a binary vector $b$ with exactly one one in each pair
by deleting one entry from every full pair and adding one to every
empty pair. Then
\[
 \|a-b\|_1=h+z=2h+m-k
 \le2SB+22SD+|m-S|.
\]
Here $|S-k|\le u$ was used. The point $V=bb^T/m$ lies in $F_m$.
To avoid division by a possibly small $S$, put $\widetilde b=(S/m)b$.
All of $r,c,\widetilde b$ have total $S$. Outer-product differences
therefore give
\begin{align*}
 \|rc^T/S-\widetilde b\widetilde b^T/S\|_1
 &\le\|r-\widetilde b\|_1+\|c-\widetilde b\|_1,\\
 \|\widetilde b\widetilde b^T/S-bb^T/m\|_1&=|S-m|.
\end{align*}
Consequently
\begin{equation}
 \|W-V\|_1\le u+v+2\|a-b\|_1+3|S-m|
 \le8mB+116mD+5|S-m|.\label{app:generator-distance}
\end{equation}
For the zero generator, any paired atom has distance $m$, and
\eqref{app:generator-distance} still holds.

Now express $W=\sum_a\lambda_aW_a$ as a finite convex combination
of generators, with totals $S_a$. By \eqref{app:exposure},
$(S_a-m)_+\le2mB(W_a)+4mD(W_a)$. Using
$|t|=2t_+-t$ gives
\[
 \sum_a\lambda_a|S_a-m|
 \le4mB(W)+8mD(W)+m-T(W).
\]
Combine their respective paired atoms in the same proportions and apply
\eqref{app:generator-distance}. The coefficients become
$8m+20m=28m$ and $116m+40m=156m$, proving
\eqref{app:face-distance}.
\end{proof}

\begin{theorem}[Outer-approximation transfer]\label{app:outer-transfer}
Let $K_{2m}\subseteq Q\subseteq K_{2m}+\eta U_1$, and let $H_m$
be the three-equation affine space defining $F_m$. For every
$Y\in Q\cap H_m$ there is $V\in F_m$ with
\[
 \|Y-V\|_1\le(184m+6)\eta.
\]
In particular, putting $A_m=m(184m+6)$ and
$P_Q=\{mY_{[m],[m]}:Y\in Q\cap H_m\}$ gives
\[
 \operatorname{COR}(m)\subseteq P_Q
 \subseteq\operatorname{COR}(m)+A_m\eta U_1.
\]
Every conic lift of $Q$ gives a lift of $P_Q$ over the same cone.
\end{theorem}
\begin{proof}
For $Y\in Q\cap H_m$, choose $W\in K_{2m}$ within distance $\eta$.
The three affine functions $B,D,T$ have coefficient norms one. Thus
$0\le B(W),D(W)\le\eta$ and $m-T(W)\le\eta$.
\Cref{app:face-rounding} supplies $V$ within $(184m+5)\eta$ of $W$,
and the triangle inequality gives the first claim. Principal-block
projection cannot increase the norm; scaling it by $m$ gives the
second inclusion. The first follows from the exact face atoms.
Adding the affine equations and composing the output map proves the
lift statement.
\end{proof}

\begin{corollary}[Uniform approximate LP size]\label{app:approx-lp}
If a polyhedron $Q$ satisfies
$K_{2m}\subseteq Q\subseteq K_{2m}+\eta U_1$ with
$\eta\le1/A_m$, every LP lift of $Q$ has $2^{\Omega(m)}$ inequalities.
\end{corollary}
\begin{proof}
The transferred polyhedron $P_Q$ lies between $\operatorname{COR}(m)$
and its entrywise unit neighborhood. For each $a\in\{0,1\}^m$ put
$M_a=2\operatorname{diag}(a)-aa^T$, whose coefficient norm is at most
one. The inequality $\langle M_a,X\rangle\le1$ is valid for the
correlation polytope: its slack at $xx^T$ equals $(1-a^Tx)^2$.
Every $X\in P_Q$ therefore satisfies $\langle M_a,X\rangle\le2$.
The fixed-factor sandwich theorem of
\citet[Theorem~6(i), author manuscript]{s6:bfps2015}, with dilation
factor two, gives $2^{\Omega(m)}$ inequalities for every LP lift of
any polyhedron in this sandwich. A lift of $Q$ transfers with no added
inequalities, proving the result.
\end{proof}

\subsection{Uniform approximate SDP size}

The previous correlation inequalities alone do not prove an SDP lower
bound. We use a different slack family and state precisely the published
quantitative input. For a nonnegative matrix $M$, its PSD rank is the
least $q$ such that $M_{st}=\langle A_s,B_t\rangle$ for PSD matrices
$A_s,B_t$ of order $q$.

The following two facts are from
\citet[Theorem~5.3 and equation~(3.11), following Theorem~3.8]
{s6:lrs-full}. For odd $k\ge3$, put
\[
 f_k(x)=\frac{(\sum_i x_i-k/2)^2-1/4}{k^2},\qquad x\in\{0,1\}^k.
\]
There is a degree-$k$ pseudo-density $D$ with
\[
 \mathbb ED=1,\quad\|D\|_\infty\le k^{3/2},\quad
 \mathbb E Df_k=-1/(4k^2),\quad
 \mathbb E Dp^2\ge0\quad(\deg p\le k/2).
\]
Expectations are uniform on the Boolean cube; a pseudo-density may be
negative. More generally, if $f:\{0,1\}^k\to[0,1]$ and a degree-$d$
pseudo-density obeys $\mathbb EDf<-\delta$ for $0<\delta\le1$,
then, for $m\ge2k$,
the matrix $M(S,x)=f(x_S)$ indexed by $k$-subsets of $[m]$ and
$x\in\{0,1\}^m$ satisfies
\begin{equation}
 \rank_{\rm psd}M\ge
 \left(\frac{c_0\delta m}{dk^2\|D\|_\infty\log m}\right)^{d/4}
 \left(\frac{\delta}{\|D\|_\infty}\right)^{3/2}
 \sqrt{\mathbb Ef},\label{app:lrs-input}
\end{equation}
where $c_0>0$ is universal. We use these established theorems without
reproving their pseudo-density machinery; all remaining transfer steps
are given below.

\begin{lemma}[Factorization of valid slacks]\label{app:slack-factor}
If a nonempty convex set $R$ has a PSD lift of order $q$, every finite
matrix $M_{st}=\ell_s(X_t)$ formed from valid nonnegative affine
functions $\ell_s$ on $R$ and points $X_t\in R$ has PSD rank at most
$q+1$.
\end{lemma}
\begin{proof}
Restrict the lifting PSD cone to its smallest face containing all feasible
matrices. That face is a PSD cone on a subspace of order $r\le q$.
A finite convex combination whose ranges span all feasible ranges is
positive definite on this subspace, giving strict feasibility there.
For each $s$, minimizing $\ell_s$ over the restricted lift has finite
nonnegative infimum. Strict primal feasibility gives dual attainment
and no duality gap. Thus, on the affine lifting space,
\[
 \ell_s(\pi(Y))=\langle B_s,Y\rangle+\tau_s,
 \qquad B_s\succeq0,\quad\tau_s\ge0.
\]
Choose a feasible lift $Y_t$ of each $X_t$. The PSD matrices
$\operatorname{diag}(B_s,\tau_s)$ and
$\operatorname{diag}(Y_t,1)$ factor $M$ with order at most $q+1$.
Facial reduction is an existence argument here; no computational
regularity assumption is required.
\end{proof}

\begin{lemma}[Robust sign-correlation lower bound]\label{app:sign-correlation}
Put $\operatorname{SCOR}(m)=\conv\{yy^T:y\in\{-1,1\}^m\}$.
If $\operatorname{SCOR}(m)\subseteq R\subseteq
\operatorname{SCOR}(m)+\rho U_1$ for $\rho\le1/2$, every PSD lift
of $R$ has order $\exp(\Omega((m/\log m)^{2/13}))$ for sufficiently
large $m$.
\end{lemma}
\begin{proof}
Fix odd $3\le k\le m/2$. For a $k$-subset $S$ define
\[
 L_S(Y)=\frac{\sum_{i,j\in S}Y_{ij}-1}{4k^2}.
\]
At $Y=(2x-\mathbf1)(2x-\mathbf1)^T$, this is $f_k(x_S)$.
Oddness of $k$ makes it nonnegative on $\operatorname{SCOR}(m)$.
Its coefficient norm is $1/(4k^2)$, so
$\ell_S=L_S+\theta$, with $\theta=\rho/(4k^2)$, is nonnegative on
$R$. By \cref{app:slack-factor}, the matrix
$M(S,x)=f_k(x_S)+\theta$ has PSD rank at most $q+1$ for a lift of
order $q$. Repeated sign-correlation points merely repeat columns.

Since $\theta\le1/(8k^2)$, the pseudo-density above gives
$\mathbb ED(f_k+\theta)\le-1/(8k^2)$.
The shifted function takes values in $[0,1]$, and its mean is at least
$(k-1)/(4k^2)\ge1/(6k)$. Apply \eqref{app:lrs-input} with
$d=k$ and $\delta=1/(16k^2)$, which satisfies the strict hypothesis.
There are universal positive constants $c_1,c_2$ such that
\[
 q+1\ge c_1k^{-23/4}
 \left(\frac{c_2m}{k^{13/2}\log m}\right)^{k/4}.
\]
Choose odd $k$ comparable to $a(m/\log m)^{2/13}$ for a sufficiently
small constant $a>0$. For all sufficiently large $m$, this is an
admissible $k$ and the parenthesized base is at least sixteen.
The resulting lower bound is $\exp(\Omega(k))$, even after the
polynomial prefactor and additive one are absorbed.
\end{proof}

\begin{theorem}[Uniform approximate SDP size]\label{app:approx-sdp}
For sufficiently large $m$, if a convex set $Q$ satisfies
$K_{2m}\subseteq Q\subseteq K_{2m}+\eta U_1$ with
$\eta\le1/(2A_m)$, every exact PSD lift of $Q$ has total matrix
order at least $\exp(\Omega((m/\log m)^{2/13}))$.
\end{theorem}
\begin{proof}
Partition matrices into four $m\times m$ blocks and use
\[
 J(W)=m(W_{11}-W_{12}-W_{21}+W_{22}).
\]
On a face atom indexed by $x$, this is
$(2x-\mathbf1)(2x-\mathbf1)^T$, so $J(F_m)=\operatorname{SCOR}(m)$.
Each input entry contributes to one output entry with coefficient
$m$ or $-m$, giving $\|J(U)-J(V)\|_1\le m\|U-V\|_1$.
The whole-matrix estimate in \cref{app:outer-transfer} consequently
shows that $R=J(Q\cap H_m)$ lies between $\operatorname{SCOR}(m)$
and its $A_m\eta\le1/2$ neighborhood. Section and projection preserve
the PSD cone, and \cref{app:sign-correlation} proves the bound.
\end{proof}

For a convex outer set, the inclusion $Q\subseteq K_{2m}+\eta U_1$
is equivalent to a uniform support gap at most $\eta$ for all cost
matrices with coefficient norm at most one. One direction follows from
the dual norm inequality; the other follows by separating a point
outside the compact convex set $K_{2m}+\eta U_1$. These are therefore
uniform objective guarantees at the displayed inverse-quadratic
accuracy. They do not exclude objective-specific relaxations, restricted
cost families, or larger approximation errors. The SDP result also
does not assert an exponential-in-$m$ lower bound for approximate
second-order-cone formulations.
